The MachRegister interface is used for accessing registers in
ProcControlAPI. The entire definition of MachRegister contains more register
names than are listed here; this appendix only lists the registers that can be
accessed through ProcControlAPI. 

An instance of MachRegister is defined for each register ProcControlAPI can
name. These instances live inside a namespace that represents the
register\'s architecture. For example, we can name a
register from an AMD64 machine with Dyninst::x86\_64::rax or a register from a
Power machine with Dyninst::ppc32::r1.


\begin{apient}
typedef unsigned long MachRegisterVal
\end{apient}

\apidesc{
The MachRegisterVal type is used to represent the contents of a register. If a
register\'s contents are smaller than MachRegisterVal, then
it will be up cast into a MachRegisterVal.
}

\begin{apient}
typedef enum { Arch_none, Arch_x86, Arch_x86_64, Arch_ppc32, Arch_ppc64 } Architecture
\end{apient}

\apidesc{
The Architecture enum represents a system\'s architecture.
}

\begin{apient}
unsigned int getArchAddressWidth(Architecture arch)
\end{apient}

\apidesc{
Returns the size of a pointer, in bytes, on the given architecture, arch.
}

\begin{apient}
MachRegister getPC(Dyninst::Architecture arch)
\end{apient}

\apidesc{
Returns the program counter (PC) for the architecture \code{arch}.
}

\begin{apient}
MachRegister getFramePointer(Dyninst::Architecture arch)
\end{apient}

\apidesc{
Returns the frame pointer for the architecture \code{arch}.
}

\begin{apient}
MachRegister getStackPointer(Dyninst::Architecture arch)
\end{apient}

\apidesc{
Returns the stack pointer for the architecture \code{arch}.
}

\begin{apient}
MachRegister getBaseRegister() const
\end{apient}

\apidesc{
This function returns the largest register that may alias with the given
register. For example, getBaseRegister on x86\_64::ah returns x86\_64::rax.
}

\begin{apient}
Architecture getArchitecture() const
\end{apient}

\apidesc{
This function returns the Architecture for this register.
}

\begin{apient}
bool isValid() const
\end{apient}

\apidesc{
This function returns true if this register is valid. Some API functions may
return invalid registers upon error.
}

\begin{apient}
MachRegisterVal getSubRegVal(MachRegister subreg,  MachRegisterVal orig)
\end{apient}

\apidesc{
Given a value for this register, orig, and a smaller aliased register, subreg,
then this function returns the value of the aliased register. For example, if
this function were called on x86::eax with subreg as x86::al and an orig value
of 0x11223344, then it would return 0x44.
}

\begin{apient}
const char *name() const
\end{apient}

\apidesc{
This function returns a human readable name that identifies this register.
}

\begin{apient}
unsigned int size() const
\end{apient}

\apidesc{
This function returns the size of the register in bytes.
}

\begin{apient}
signed int val() const
\end{apient}

\apidesc{
This function returns a unique integer that represents this register. This can
be useful for writing switch statements that take a MachRegister. The unique
integers for a MachRegister can be found by preceding the register object name
with an ``i''. For example, a switch statement for MachRegister, reg, could be written as:
}


\begin{lstlisting}[language=c++]
switch (reg.val()) {
  case x86\_64::irax:
  case x86\_64::irbx:
  case x86\_64::ircx:
   ...
}
\end{lstlisting}

\begin{apient}
bool isPC() const
\end{apient}

\apidesc{
Returns \code{true} if the register is the program counter (PC) for any architecture.
}

\begin{apient}
bool isFramePointer() const
\end{apient}

\apidesc{
Returns \code{true} if the register is the frame pointer for any architecture.
}

\begin{apient}
bool isStackPointer() const
\end{apient}

\apidesc{
Returns \code{true} if the register is the stack pointer for any architecture.
}
